% !TeX root = main.tex
% =====================================================================
%  Präambel – Pakete und Layout (wie Bachelorarbeit)
% =====================================================================

% ---------- Sprache & Schrift ----------
\usepackage[utf8]{inputenc}
\usepackage[T1]{fontenc}
\usepackage[ngerman]{babel}
\usepackage{csquotes}                 % deutsche Anführungszeichen (auch für biblatex)

% ---------- Seitenlayout nach Knechtel-Vorlage ----------
\usepackage{geometry}
\geometry{left=3cm,right=2.5cm,top=2.5cm,bottom=2.0cm}
\usepackage{setspace}
\setstretch{1.5}

\usepackage{fancyhdr}
\pagestyle{fancy}
\fancyhf{}
\fancyfoot[C]{\thepage}
% Dünne Linie oben wie in der BA. Entfernen mit:
% \renewcommand{\headrulewidth}{0pt}

% ---------- Inhaltsverzeichnis ----------
\usepackage{tocloft}
\setlength{\cftbeforesecskip}{6pt}

% ---------- Fußnoten ----------
\usepackage[bottom]{footmisc}
\interfootnotelinepenalty=10000

% ---------- Tabellen ----------
\usepackage[table]{xcolor}            % lädt auch colortbl
\usepackage{array}
\usepackage{tabularx}
\usepackage{booktabs}
\newcolumntype{L}[1]{>{\raggedright\arraybackslash}p{#1}}
\newcolumntype{C}[1]{>{\centering\arraybackslash}p{#1}}

% ---------- Abbildungen ----------
\usepackage{graphicx}
\graphicspath{{./abbildungen/}}
\usepackage{float}
\usepackage{caption}
\usepackage{subcaption}
\usepackage{adjustbox}
\usepackage[percent]{overpic}
\usepackage{pdfpages}
\usepackage{pdflscape}
\usepackage{pgfplots}
\pgfplotsset{compat=1.18}

% ---------- Mathematik & Einheiten ----------
\usepackage{amsmath,amssymb}
\usepackage{siunitx}
\sisetup{locale=DE, per-mode=symbol}
\DeclareSIUnit{\kilogramforce}{kgf}

% ---------- Listen ----------
\usepackage{enumitem}
\setlist[itemize]{noitemsep, topsep=0pt}

% ---------- Sonstiges ----------
\usepackage{comment}

% Silbentrennung für Problemwörter, z. B.:
% \hyphenation{Wieder-holbar-keits}

% ---------- Quellcode ----------
\usepackage{listings}
\lstset{
	language=Python,
	basicstyle=\ttfamily\small,
	keywordstyle=\color{blue},
	commentstyle=\color{gray},
	stringstyle=\color{green!50!black},
	numbers=left,
	numberstyle=\tiny,
	breaklines=true,
	frame=single
}

% ---------- Literatur ----------
\usepackage[backend=biber]{biblatex}
\addbibresource{literatur.bib}

% ---------- Hyperref (möglichst zuletzt laden) ----------
\usepackage{hyperref}
% Für eine Druckversion ohne farbige Linkrahmen: \hypersetup{hidelinks}
\hypersetup{
	pdftitle={\ArbeitTitel},
	pdfauthor={\AutorName}
}
